% Replacement for sections/motivation.tex in the v0.11-I paper source.
% Text revision only. This file has not been compiled in this review session.
% Author TODOs are visible by default; detailed closure criteria are in TODO_REGISTER.md.
\providecommand{\dtctodo}[2]{\textbf{[TODO #1: #2]}}
\section{Allocation and Retention Trade-off}
\label{sec:motivation}

\textit{Waiting capacity and cache capacity.}
The BICG controls in Section~IV separate two resource choices~\cite{dtc-artifact}.
Increasing the conventional PIB from 8 to 128 entries retains the 16-KiB cache and 32 MSHRs.
PIB-full events disappear, yet replay takes 5.43\% more cycles.
A separate 80-KiB cache gives a $3.002\times$ speedup over the 16-KiB cache.
It retains eight PIB entries and 32 MSHRs.
PIB capacity alone does not explain the performance difference.
% TODO CMP-01: These are accepted resource contrasts, not a hit-rate-only attribution.
% Preserve the nonbinding cap and inactive-option qualifications in the existing methods.
\dtctodo{CMP-01}{Explain the gap to conventional designs}

\textit{Lookup and retained data.}
Suppose two admitted instructions share a physical line.
Each PIB entry records its index and dependencies, not a private copy of the data~\cite{zhao2024}.
If the line's tag is replaced, both instructions keep their physical references.
Later lookups cannot find that allocation.
The evaluated DTC read path has 640 physical lines but only 128 tag mappings.
Its pool can retain 80~KiB, but at most 16~KiB can remain mapped.
Pool size and lookup coverage therefore need separate evaluation.
% This is a worked structural example, not a measured request trace or a proof of savings.
% TODO COV-01: Consume the complete accepted T1 matrix before claiming a useful Tag benefit.
% TODO HW-01: Include indices, dependencies, waiter state, ports and the ordinary side array.
\dtctodo{COV-01}{Tag-coverage evidence}
\dtctodo{HW-01}{Storage and service cost}

\textit{Readiness and release.}
The IO/OO BICG control fixes the PIB at 128 entries and the pool at 640 lines~\cite{dtc-artifact}.
Mean ready-to-collection wait is 3,557.95 core cycles for IO and 1.16 for OO.
These means cover retired cacheable warp instructions.
Each interval starts after all references have been admitted and their data are ready.
The instruction still occupies its PIB entry during this wait.
IO can also retain lines assigned to its release list.
Completion and reclamation therefore require evaluation at matched waiting depth.
Figure~\ref{fig:design} illustrates the different release points.
% TODO ORG-01: BICG supports the matched-depth observation; GEMM T0 remains pending review.
% A short post-ready interval does not identify the full issue-to-completion latency.
% This comparison does not by itself close CMP-01, prove net area saving or establish generality.
\dtctodo{ORG-01}{GEMM matched-depth evidence}
